\section{A direct bound from the critical source norm}
\label{r:app:critical-source}

The source can also be tested before velocity-intersection membership is
known. Begin with a classical forced history on its interval of existence
and retain the nonlinear current in the critical energy balance. Put
\[
 y=\norm u_{\dot H^{1/2}},\qquad
 z=\norm u_{\dot H^{3/2}},\qquad
 \beta=\norm{\Pp f}_{\dot H^{-1/2}},\qquad
 C_*=\frac1{\sqrt2\pi}.
\]
The product estimate used in \eqref{r:eq:critical-production}, applied
before substituting any rectangle bound, gives
\begin{equation}
 y y'+(\nu-C_*y)z^2\le\beta z.
 \label{r:eq:critical-source-energy}
\end{equation}
This calculation is on a strict classical interval. It assumes no bound
on a velocity coordinate at the terminal time.

As long as \(y<\nu/C_*\), completing the square in \(z\) yields
\[
 y y'\le\frac{\beta^2}{4(\nu-C_*y)}.
\]
Multiplication by the positive denominator makes the left side an exact
derivative. Define
\[
 \Psi(s)=2\nu s^2-\frac43C_*s^3,\qquad
 B(t)=\int_0^t\beta(s)^2\dd s.
\]
Then
\begin{equation}
 \Psi(y(t))\le\Psi(y(0))+B(t),\qquad
 \Psi'(s)=4(\nu-C_*s)s.
 \label{r:eq:critical-source-primitive}
\end{equation}
The function \(\Psi\) is strictly increasing between zero and
\(\nu/C_*\). At the upper endpoint it takes the value
\[
 \Psi(\nu/C_*)=\frac{2\nu^3}{3C_*^2}
 =\frac{4\pi^2}{3}\nu^3.
\]

\begin{theorem}[Critical source criterion]\label{r:thm:critical-source}
Under the datum and source assumptions of this paper, suppose
\(y_0=\norm{u_0}_{\dot H^{1/2}}<\nu/C_*\) and
\begin{equation}
 B(T)<\frac{4\pi^2}{3}\nu^3-\Psi(y_0).
 \label{r:eq:critical-source-condition}
\end{equation}
The forced history extends smoothly through \(T\) and belongs to
\(\mathscr I_T\). In particular, for rest data the sufficient condition is
\(B(T)<4\pi^2\nu^3/3\).
\end{theorem}

\begin{proof}
Let \(r<\nu/C_*\) be the unique nonnegative number satisfying
\(\Psi(r)=\Psi(y_0)+B(T)\). If \(y\) first exceeded \(r\) before
reaching \(\nu/C_*\), \eqref{r:eq:critical-source-primitive} would fail.
Continuity thus gives \(y\le r\) throughout classical existence up
to \(T\). At a zero of \(y\), the energy identity for \(y^2\) and
absolute continuity justify the integrated inequality without dividing
by \(y\).

Write \(a=\nu-C_*r>0\). Young's inequality in
\eqref{r:eq:critical-source-energy} now gives
\[
 y y'+\frac a2 z^2\le\frac{\beta^2}{2a},\qquad
 \int_0^S z^2\dd t\le\frac{y_0^2}{a}+\frac{B(T)}{a^2}
 \quad(S<T).
\]
These are bounds obtained from the source and initial datum.
To see directly how they prevent an endpoint loss, put
\(Y=\norm{\nabla u}_2\) and \(D=\norm{\Delta u}_2\).
Differentiation of the momentum equation, integration by parts, and
incompressibility give
\[
 \frac12(Y^2)'+\nu D^2
 \le\norm{\nabla u}_3\norm{\nabla u}_2
       \norm{\nabla u}_6+\norm{\Pp f}_2D.
\]
Let \(C_S\) be a finite constant in
\(\norm{\nabla u}_6\le C_S\norm{\Delta u}_2\), and let
\(C_3=2^{-1/6}\pi^{-1/3}\) as above. Since
\(\norm{\nabla u}_3\le C_3z\), two applications of Young's inequality
give
\[
 (Y^2)'+\nu D^2
 \le\frac{2(C_SC_3)^2}{\nu}z^2Y^2
       +\frac2\nu\norm{\Pp f}_2^2.
\]
Gronwall's inequality bounds \(Y\) uniformly. The ordinary energy
identity also gives
\(\norm{u(t)}_2\le\norm{u_0}_2+
\int_0^t\norm{\Pp f(s)}_2\dd s\).
Thus the full \(H^1\) norm remains bounded.

For completeness, this bound supplies a uniform local continuation time.
The estimate
\(\norm{(v\cdot\nabla)w}_{3/2}
\le C\norm v_{H^1}\norm w_{H^1}\), the boundedness of \(\Pp\)
on \(L^{3/2}\), and the heat estimates from \(L^{3/2}\) to \(H^1\)
bound the bilinear mild map on a time interval of length \(h\) by
\[
 C_\nu(h^{1/4}+h^{3/4})
 \norm v_{C_tH^1}\norm w_{C_tH^1}.
\]
The two powers integrate the heat factors \(t^{-3/4}\) and
\(t^{-1/4}\). A contraction ball of fixed radius now has a fixed
positive existence time, using also the bounded \(H^1\) norm of the
source. Restarting near a proposed finite endpoint continues the
history beyond it. Heat smoothing at positive elapsed time and the
smooth source preserve every finite spatial order; differentiation
of the equation then preserves the full temporal jet. Together with
the smooth initial datum and source, this proves \(\mathscr I_T\)
membership and the endpoint conclusions.
\end{proof}

The threshold is unchanged by the Navier--Stokes rescaling
\(u_\lambda(x,t)=\lambda u(\lambda x,\lambda^2t)\),
\(f_\lambda(x,t)=\lambda^3f(\lambda x,\lambda^2t)\).
The squared \(\dot H^{-1/2}\) source norm gains a factor
\(\lambda^2\), which its time integral cancels.

\subsection{A necessary \texorpdfstring{\(H^3\)}{H3} source budget}

For a candidate that starts from rest, write
\[
 Y(t)=\norm{u(t)}_{H^3},\qquad
 D(t)=\norm{\nabla u(t)}_{H^3},\qquad
 g(t)=\norm{\Pp f(t)}_{H^3},\qquad
 G_T=\int_0^Tg(t)\dd t.
\]
On each strict classical interval, the momentum equation and
\eqref{r:eq:product} give
\[
 YY'+\nu D^2\le8c_\infty Y^2D+gY,
 \qquad c_\infty=(8\pi)^{-1/2}.
\]
Completing the entire viscous square yields
\[
 8c_\infty Y^2D-\nu D^2
 \le\frac{16c_\infty^2}{\nu}Y^4
 =\frac2{\pi\nu}Y^4.
\]
Thus \(Y'\le2Y^3/(\pi\nu)+g\) wherever \(Y>0\).
At zeros, use \(Y_\varepsilon=(Y^2+\varepsilon^2)^{1/2}\);
its derivative obeys
\(Y_\varepsilon'\le2Y_\varepsilon^3/(\pi\nu)+g\).

If \(G_T=\infty\), the finite source condition already fails. Otherwise
put
\[
 Z_\varepsilon(t)=Y_\varepsilon(t)+\int_t^Tg(s)\dd s.
\]
Since \(Y_\varepsilon\le Z_\varepsilon\),
\[
 Z_\varepsilon'\le\frac2{\pi\nu}Z_\varepsilon^3,
 \qquad Z_\varepsilon(0)=G_T+\varepsilon.
\]
Integrating this scalar comparison and taking \(\varepsilon\downarrow0\)
gives
\begin{equation}
 Y(t)\le
 \frac{G_T}{\sqrt{1-4G_T^2t/(\pi\nu)}}
 \quad\text{if}\quad
 G_T<\sqrt{\frac{\pi\nu}{4T}}.
 \label{r:eq:direct-H3-budget}
\end{equation}
This bound uses the finite source budget and strict-interval equations;
it assumes neither a terminal velocity norm nor smooth continuation of
force through \(T\).

The calculated exterior has an unbounded \(L^3\) norm as \(t\uparrow1\).
Since \(H^3\) controls \(\dot H^{1/2}\), which embeds into \(L^3\),
its \(H^3\) norm is unbounded as well. Equation~\eqref{r:eq:direct-H3-budget}
thus proves the following lower bound for its actual residual:
\begin{equation}
 G_{\rm OA}\ge\frac{\sqrt{\pi\nu}}2.
 \label{r:eq:actual-H3-budget-failure}
\end{equation}

\subsection{Strict failure of the critical source condition}

For the prescribed exterior, cylindrical integration on
\(\mathcal E_\tau\) gives an explicit critical lower bound:
\[
 \int_{\mathcal E_\tau}|u|^3\dd x
 =C_L\tau^{-4h},\qquad
 C_L=4\pi d\int_c^{2c}k(y)^3y\dd y>0.
\]
The Sobolev inequality implies
\[
 \norm u_{\dot H^{1/2}}
 \ge C_3^{-1}C_L^{1/3}\tau^{-4h/3}\longrightarrow\infty.
\]
The viscosity rescaling preserves this divergence. Continuity gives a first
time \(S<1\) at which
\(y(S)=\nu/C_*\), with \(y(t)<\nu/C_*\) for \(0\le t<S\).
Integrate \eqref{r:eq:critical-source-primitive} before \(S\), then pass to
\(S\) by continuity. The rest datum gives
\begin{equation}
 B(S)\ge\Psi(\nu/C_*)=\frac{4\pi^2\nu^3}{3}.
 \label{r:eq:necessary-source-budget}
\end{equation}
The bound follows from the explicit exterior and momentum equation on
strict subintervals of \([0,1)\).

The residual continues to contribute a positive source norm after this
crossing. Use the divergence-free test field \(w\) from the fixed
exterior calculation leading to \eqref{r:eq:actual-pressure-obstruction},
and write \(\gamma_w:=\lim_{t\uparrow1}\ip{f(t)}w>0\).
For some \(\eta_1>0\), local convergence on its support gives
\(\ip{f(t)}w\ge\gamma_w/2\) for \(1-\eta_1<t<1\), independently
of terminal convergence elsewhere. Homogeneous Sobolev duality gives
\[
 \frac{\gamma_w}{2}\le\ip{\Pp f(t)}w
 \le\norm{\Pp f(t)}_{\dot H^{-1/2}}\norm w_{\dot H^{1/2}}.
\]
The test norm is finite and positive because \(w\) is smooth, compactly
supported, and nonzero. Let \(a=\max\{S,1-\eta_1\}<1\).
Integrating over \((a,1)\), which is disjoint from \((0,S)\), yields
\begin{equation}
 \begin{aligned}
 B(1)&\ge\frac{4\pi^2\nu^3}{3}
       +\frac{(1-a)\gamma_w^2}{4\norm w_{\dot H^{1/2}}^2}\\
     &>\frac{4\pi^2\nu^3}{3}.
 \end{aligned}
 \label{r:eq:strict-critical-budget-failure}
\end{equation}
For general viscosity, use the rescaled exterior and test field; the
source budget and positive correction both scale by \(\nu^3\).
The calculation also permits \(B(1)=\infty\). Thus no global terminal
source regularity is needed to establish failure of this sufficient
criterion. The source may be smooth and have finite norms, but these
particular quantitative estimates do not cover its magnitude. This failure
does not restrict the global theorem of Part~II.
